\section{All closed faces from two elementary matrix lifts}
\label{sec:closed-faces-short}

The noncoercive faces do not require a list of support masks. The
physical classification in Proposition~\ref{prop:compat} leaves just
two kinds of omitted exclusive pairs, and both admit short exact lifts.

\begin{lemma}[Two kinds of exclusive pair]
\label{lem:closed-face-orbits}
Every pair of exclusive losing events is CHSH-equivalent either to
$\{\lambda,\alpha\}$ or to $\{\lambda,m\}$. Consequently every
noncoercive source belongs to a relabeling of one of these two closed
coefficient faces.
\end{lemma}
\begin{proof}
If the two events have the same input pair, they are the two opposite
losing outcomes at that input. There are four such pairs. Applying
$g_A,g_B$ from Section~\ref{sec:physical} to
$\{\lambda,\alpha\}$ gives $\{m,D\}$ and $\{n,A\}$; applying
$g_A$ after $g_B$ gives $\{c,\beta\}$.

Otherwise exclusivity means opposite outcomes at one shared local
measurement and different remote settings. Each of the four local
measurements gives two such pairs. Starting from $\{\lambda,m\}$,
the words
\[
 1,\quad B,\quad T,\quad BA,\quad BT,\quad TA,\quad TB,\quad BAB
\]
in $g_A,g_B,g_T$, applied from left to right, give respectively
\[
 \{\lambda,m\},\ \{A,\beta\},\ \{c,m\},\ \{c,n\},\
 \{\alpha,n\},\ \{A,\lambda\},\ \{D,\beta\},\ \{D,\alpha\}.
\]
These are all eight possibilities. Every noncoercive omission set
contains an exclusive pair by Proposition~\ref{prop:compat}; deleting
further coefficients only takes a closed subface.
\end{proof}

We use the following elementary matrix fact. Its proof is included to
avoid importing a matrix factorization theorem.

\begin{lemma}[An exact quadratic matrix lift]
\label{lem:quadratic-matrix-lift}
Let $K(t)=K_0+tK_1+t^2K_2$ be a real symmetric $r$ by $r$ matrix
polynomial. Then $K(t)\succeq0$ for every real $t$ if and only if
there is a skew-symmetric $S$ such that
\[
 \begin{pmatrix}
 K_0&K_1/2+S\\ K_1/2-S&K_2
 \end{pmatrix}\succeq0.                         \tag{CF1}
\]
\end{lemma}
\begin{proof}
The reverse direction is obtained by evaluation on $(z,tz)$.
For the forward direction let $\mathcal C$ be the cone of coefficient
triples produced by positive semidefinite Gram matrices. This cone is
closed: for a convergent sequence of coefficient triples the Gram
trace equals $\operatorname{tr}(K_0+K_2)$, so a sequence of realizing
positive semidefinite Grams is bounded and has a convergent subsequence.

A linear functional nonnegative on $\mathcal C$ has the form
\[
 \ell(K)=\operatorname{tr}(U K_0)+\operatorname{tr}(V K_1)
                         +\operatorname{tr}(W K_2),\qquad
 H=\begin{pmatrix}U&V\\V&W\end{pmatrix}\succeq0,
                                                        \tag{CF2}
\]
where $U,V,W$ are symmetric. Positivity of $H$ implies
$\operatorname{ran}V\subseteq\operatorname{ran}U$ and
$W-VU^+V\succeq0$, where $U^+$ is its inverse on its range and
zero on its kernel. On $\operatorname{ran}U$ diagonalize the symmetric
matrix $U^{+/2}VU^{+/2}$. If $(t_j,e_j)$ are its orthonormal
eigenpairs and $z_j=U^{1/2}e_j$, then
\[
 \begin{pmatrix}U&V\\V&VU^+V\end{pmatrix}
   =\sum_j\binom{z_j}{t_jz_j}\binom{z_j}{t_jz_j}^{\!T}.       \tag{CF3}
\]
The remaining block $W-VU^+V$ is a sum of $w_iw_i^T$. Thus
\[
 \ell(K)=\sum_j z_j^TK(t_j)z_j+\sum_iw_i^TK_2w_i\ge0,
\]
because $K_2\succeq0$ follows by taking the leading coefficient at
infinity. Finite-dimensional separation of the closed convex cone
$\mathcal C$ now gives membership in $\mathcal C$, which is (CF1).
This proof also handles $U=0$ and singular $U$ literally.
\end{proof}

For each fixed matrix size the lemma is a first-order statement in
ordered-field arithmetic. It therefore holds over every real closed
coefficient field as well. All changes of variables used below are
within that one field; no inverse in a tensor product is introduced.

\begin{theorem}[The missing closed face has a six-dimensional lift]
\label{thm:closed-face-adjacent}
On the closed canonical face $\lambda=m=0$, with every remaining
event coefficient and $d$ nonnegative, physicality is equivalent to
(CF1) for $r=3$ and
\begin{align*}
 K_0&=\operatorname{diag}(4c+4D+2d,\ 4A,\ 4\alpha),\\
 K_2&=\operatorname{diag}(4c,\ 4A+4\alpha+2d,\ 4D),\\
 K_1&=\begin{pmatrix}
 0&0&-(4D+2d)\\
 0&0&-(4\alpha+2d)\\
 -(4D+2d)&-(4\alpha+2d)&0
 \end{pmatrix}.                                           \tag{CF4}
\end{align*}
There are only three auxiliary variables, the independent entries of
$S$. In particular this is an affine semidefinite lift, with no
nonlinear inequalities in the source coefficients.
\end{theorem}
\begin{proof}
The case $d=0$ is immediate from the original event-square sum and
from the positive-square expression below, so first suppose $d>0$.
In the parent set $W=a/u$ and send $u$ to zero while holding
$a,X,Y,Z,v$ fixed. Put $R=a-Y$. The limiting form is a nonnegative
multiple $2t_5(1+v^2)X^2$, with
$t_5=d+2(\beta+n)>0$, plus
\[
 \begin{split}
 G_v(R,Y,Z)={}&4c(1+v^2)R^2+4A(1+v^2)Y^2
                  +(4D+2d)(R-vZ)^2\\
 &+(4\alpha+2d)(vY-Z)^2-2d(1+v^2)Z^2.
 \end{split}                                             \tag{CF5}
\]
Its matrix is precisely $K(v)$ from (CF4). Homogeneity of the state
makes this limiting test legitimate. Hence physicality implies
$K(v)\succeq0$ for every real $v$.

Conversely, in this face the Laurent determinant is
\[
 L(U,V)=h+q_{10}/V+q_{12}V
                   +U(q_{20}/V+q_{21}+q_{22}V),             \tag{CF6}
\]
and every displayed $q$ is nonnegative. The same boundary limit in
the determinant identity of Section~\ref{sec:physical} gives exactly
\[
 \det K(v)=\frac{8(1+v^2)}{t_5}
                (q_{10}+hv^2+q_{12}v^4).                  \tag{CF7}
\]
Indeed the transformed parent matrix has determinant
$\det P/u^2$, and its limit is
$2t_5(1+v^2)\det K(v)$. Thus (CF7) requires no expansion of a
new determinant. If $K(v)\succeq0$, (CF7) makes the first three
terms of (CF6) nonnegative for $V>0$; all the terms multiplied by
$U$ are nonnegative. The determinant criterion, whose three hard
eigenvalues are positive for $d>0$ and $uv\ne0$, proves positivity
of the full parent. Polynomial continuity supplies every finite or
projective angle endpoint. Lemma~\ref{lem:quadratic-matrix-lift}
now proves the asserted lift.

For clarity, no strictness or nonzero event coefficient was needed.
When $d=0$, (CF5) is itself a sum of squares linear in $v$, so it
directly supplies a Gram in (CF1). Literal zeros are therefore
included without a limiting source argument.
\end{proof}

The lift also has a strictly feasible coefficient point. Set
$c=A=2$, $D=\alpha=d=n=\beta=1$, and $S=0$. A permutation of
its six coordinates gives two three-dimensional blocks. In each,
the two diagonal hard pivots are $14,8$ and their Schur complement
is
\[
 4-\frac9{14}-\frac98=\frac{125}{56}>0.                    \tag{CF8}
\]
Thus the usual attained polar-duality argument applies to this lift;
it does not require an approximate or numerical certificate.

\begin{theorem}[The opposite-outcome face has two three-dimensional blocks]
\label{thm:closed-face-opposite}
On the closed canonical face $\lambda=\alpha=0$, define
\[
 p=4c+4D+2d,\qquad q=4m+4\beta+2d,\qquad r=4D+2d.
\]
Physicality is equivalent to the existence of one real $\zeta$ with
\[
 \begin{pmatrix}
 4m&-4m&\zeta\\
 -4m&q+4n&q\\
 \zeta&q&q
 \end{pmatrix}\succeq0,\qquad
 \begin{pmatrix}
 p+4A&4D&-p\\
 4D&4D&-(4m+r)-\zeta\\
 -p&-(4m+r)-\zeta&p
 \end{pmatrix}\succeq0.                                  \tag{CF9}
\]
All original coefficient signs are retained.
\end{theorem}
\begin{proof}
Set $u=\epsilon s$, $v=\epsilon$ and
$(W,Z)=(w,z)/\epsilon$, while $X=x,Y=y$. The finite limit of
the parent as $\epsilon$ tends to zero is
\[
\begin{split}
G_s={}&4nx^2+4m(w-x-sz)^2+4c(sw-y)^2+4Ay^2\\
 &+r(sw-y-z)^2+(4\beta+2d)(x+sz)^2-4dyz-2dz^2.
\end{split}                                               \tag{CF10}
\]
The divergent terms cancel between the positive signed packet and
the negative core. Write $J(s)$ for the matrix of $G_s$ on
$(w,x,y,z)$. Its leading coefficient is
$\operatorname{diag}(p,0,0,q)$, and its determinant is
\[
 \det J(s)=16(q_{02}+hs^2+q_{20}s^4).                      \tag{CF11}
\]
Again this follows directly from the original determinant identity:
the state change multiplies the determinant by $\epsilon^{-4}$,
and only the displayed three lowest-order terms survive.

Physicality implies $J(s)\succeq0$ for every real $s$. Conversely,
if $d>0$ and $J(s)\succeq0$, (CF11) implies
$q_{02}/s^2+h+q_{20}s^2\ge0$ for $s\ne0$. The full Laurent
polynomial in this face is
\[
 h+q_{02}V/U+q_{20}U/V+q_{12}V+q_{21}U+q_{22}UV.
\]
Take $s^2=U/V$. Its first three terms are nonnegative and its
remaining coefficients are nonnegative. The determinant criterion
and continuity prove positivity of the entire parent. If $d=0$,
the parent and (CF10) are event-square sums, so the same equivalence
is immediate.

Apply Lemma~\ref{lem:quadratic-matrix-lift} to $J(s)$. Its Gram
initially acts on $(w,x,y,z,sw,sx,sy,sz)$. The two diagonal entries
for $sx,sy$ are zero; positivity makes their entire rows and columns
zero. The coefficient equations leave only one free off-diagonal
entry, $\zeta$. Permuting the remaining coordinates into
$(w,x,sz)$ and $(y,z,sw)$ gives exactly the two blocks (CF9).
Equivalently, direct multiplication of the displayed blocks gives
the exact identity
\[
 G_s=(w,x,sz)B_1(w,x,sz)^T+(y,z,sw)B_2(y,z,sw)^T,
\]
where $B_1,B_2$ are those blocks. The terms involving $\zeta$
cancel because $w(sz)=z(sw)$. This completes necessity as well
as sufficiency of the affine lift.
\end{proof}

This lift is strictly feasible too: take
$c=A=n=\beta=2$, $m=D=d=1$, and $\zeta=-5$. The leading
principal minors of the two blocks are respectively
$(4,72,234)$ and $(22,72,234)$, all positive integers.

\begin{corollary}[The entire closed-face part is explicit]
\label{cor:all-closed-faces}
Every noncoercive canonical source belongs to a closed coefficient
face with an explicit affine semidefinite lift. The two representatives
are the two three-dimensional blocks (CF9) for $\lambda=\alpha=0$
and the single six-dimensional block (CF1)--(CF4) for $\lambda=m=0$.
All further literal event zeros are included by linear restriction.
\end{corollary}
\begin{proof}
Apply Lemma~\ref{lem:closed-face-orbits} and
Theorems~\ref{thm:closed-face-adjacent}--\ref{thm:closed-face-opposite}.
Relabeling is an invertible
linear change in the original Bell coordinates. The finite union of
these coefficient faces is used only in the canonical source
reduction; no union of spectrahedra is asserted to be convex.
\end{proof}
